\section{Cosine \texorpdfstring{$\rone$}{R1}: Structural Failure}
\label{sec:cosine_failure}

Section~\ref{sec:direction_quality} established that $\dsyc$ achieves AUROC 0.827 on Qwen3-8B and generalizes across models.
We now test whether this discriminability translates into a viable GRPO reward signal.
This section reports four configurations (F1a--F1d) that form a systematic diagnostic chain: each configuration resolves the failure mode of its predecessor but exposes a deeper structural limitation.
Table~\ref{tab:cosine_chain} summarizes the chain; Figure~\ref{fig:cosine_failure} visualizes the training dynamics.

\begin{table}[t]
\centering
\caption{Cosine $\rone$ failure chain. Each configuration addresses the preceding failure but reveals a more fundamental limitation. All experiments use Qwen3-8B with LoRA (rank$=$64) and GRPO with $R = w_1 \cdot \rone + w_3 \cdot \rthree$.}
\label{tab:cosine_chain}
\begin{tabular}{lllll}
\toprule
Config & $\rone$ Formula & $w_1$ & Outcome & Root Cause \\
\midrule
F1a & $\tanh(\mathbf{h}\cdot\dsyc/20)$ & 0.25 & Gate saturation & Norm mismatch \\
F1b & $\cos(\mathbf{h},\, \dsyc)$ & 0.25 & $\rone \approx 0$ & Gap too small \\
F1c & +recalibration & 0.25 & Mode collapse & Reward discontinuity \\
F1d & $5\cos(\mathbf{h},\, \dsyc)$ & 0.40 & Mode collapse & $\rthree$ saturation + noise \\
\bottomrule
\end{tabular}
\end{table}

\paragraph{F1a: Tanh gate saturation (exp-002).}
The initial configuration defines $\rone = \tanh(\mathbf{h} \cdot \dsyc / 20)$ with $w_1{=}0.25$, $w_3{=}0.75$.
The dot product formulation preserves directional sensitivity but is not scale-invariant: activation norms in the residual stream range from $100{\times}$ to $1{,}200{\times}$ the direction norm across layers.
This norm mismatch drives the pre-activation into the saturated tails of tanh, where the gradient vanishes.
By step 955, $\rone$ contributes effectively zero gradient to the combined reward, and training was terminated.
Subsequent weight adjustments ($w_1{=}0.15$, $w_3{=}0.65$) did not resolve the saturation because the pre-activation magnitude, not the weight, controls the gate.
The diagnostic conclusion is clear: any formulation that passes raw dot products through a bounded nonlinearity will saturate in high-norm activation spaces.
Cosine similarity, which normalizes by both norms, is the natural alternative.

\paragraph{F1b: Raw cosine signal too weak (exp-002v2).}
Replacing the dot product with cosine similarity ($\rone = \cos(\mathbf{h},\, \dsyc)$) eliminates the norm dependence.
However, the cosine gap between valid corrections and sycophantic capitulations is only 0.098 when averaged across the seven highest-AUROC layers.
With $w_1{=}0.25$, this yields an effective $\rone$ contribution of approximately $0.025$ per token, a signal that is negligible relative to $\rthree$ variance.
Throughout training, $\rone \approx 0$, and the combined reward $R$ is functionally equivalent to $w_3 \cdot \rthree$ alone.
Training was terminated after confirming that the sycophancy component exerts no measurable influence on policy updates.
The failure isolates a quantitative limitation: the cosine projection of $\dsyc$ onto individual activations produces a gap too narrow to generate meaningful reward contrast between sycophantic and non-sycophantic completions.

\paragraph{F1c: Recalibration induces mode collapse (exp002v2-recalib).}
One hypothesis for the weak signal in F1b is that $\dsyc$ drifts during LoRA training, eroding the already-small gap.
To test this, we resume from the exp-002v2 checkpoint at step 500 and enable the online recalibration procedure described in \S\ref{ssec:recalib}, re-extracting $\dsyc$ every 500 steps.
Recalibration increases the offline AUROC from 0.82 to 0.91, confirming that the refreshed direction better separates the updated activations.
Training collapses nonetheless.
Within 300 steps of the first recalibration event, the policy degenerates into a low-entropy mode: generation entropy drops to 0.069, and mean response length contracts to 62 tokens.
The mechanism is a reward discontinuity at the recalibration boundary.
When $\dsyc$ is replaced by a fresh direction, the cosine projection of every cached rollout shifts abruptly.
Completions that were mildly rewarded under the stale direction may become penalized under the fresh one, and vice versa.
This discontinuity destabilizes the GRPO advantage estimates and drives the policy toward a narrow region of activation space that remains low-risk under both the old and new directions.
The result is a degenerate policy that produces short, hedged responses regardless of the input.
Improving direction quality via recalibration, paradoxically, makes the reward landscape less stable.

\paragraph{F1d: Scaling amplifies noise (exp-002v3).}
Configuration F1d addresses the magnitude problem directly: $\rone = 5 \cdot \cos(\mathbf{h},\, \dsyc)$ with $w_1{=}0.4$, $w_3{=}0.6$.
The $5{\times}$ scaling lifts the $\rone$ contribution above the $\rthree$ noise floor, but it simultaneously amplifies the noise in the cosine projection.
Training proceeds in three phases.
In the first phase (steps 100--300), $\rthree$ saturates at 0.9998 as the policy rapidly learns to produce NLI-consistent answers; mean response length drops from 247 to 80 tokens and entropy falls from 0.45 to 0.10.
In the second phase (steps 300--2150), the combined reward degenerates to $R \approx w_1 \cdot \rone + \text{const}$ because the saturated $\rthree$ provides no gradient.
$\rone$ is now the sole learning signal, but its mean oscillates between $-0.46$ and $+0.41$ with no sustained positive trend, because the cosine projection conflates warranted belief revision with unwarranted capitulation (both shift activations toward the user's position along $\dsyc$).
In the third phase (step 2160), the policy collapses terminally: entropy drops from 0.056 to 0.026, and the fraction of completions with zero reward standard deviation rises from 0.125 to 0.419, indicating that the policy has converged to a near-deterministic mode.

\paragraph{Unified diagnosis.}
The four configurations converge on a single structural limitation.
The cosine gap between valid and sycophantic activations (0.098 across layers) is too narrow to serve as a reward signal, and this narrowness is not a calibration artifact.
Recalibration raises offline AUROC to 0.91 yet rewards still fail (F1c), confirming that direction quality is not the bottleneck.
The root cause is semantic: $\dsyc$ encodes the \emph{direction} of belief change (toward or away from the user's position) but does not encode whether that change is \emph{warranted}.
A model that appropriately revises an incorrect answer in response to strong evidence and a model that capitulates under social pressure both shift their activations toward the user's position along $\dsyc$.
The cosine projection assigns similar scores to both behaviors.
Scaling the projection (F1d) amplifies this conflation rather than resolving it.
Any reward formulation that uses $\dsyc$ as a continuous scalar will inherit this ambiguity.
Section~\ref{sec:probe_trap} examines whether discretizing the signal via a binary probe can escape this limitation.
